\section{Implementation}
\label{sec:implementation}

We provide three implementations of \textsc{Kaiburr}, each forked from an existing ML-KEM implementation: 
a C implementation, a Jasmin implementation with machine-checked constant-time verification, and an assembly 
implementation for Pavona's Asymmetric Cryptography Coprocessor (ACC).
\footnote{\url{https://github.com/pavona/pavona}}
The ACC reuses ML-KEM's 
arithmetic core and accelerates Keccak in hardware. This provides us with a complementary hardware-accelerated setting 
where we can evaluate the cost of scaling the module rank $k$. As upstream, the C and Jasmin code each have a 
portable \texttt{ref} path and an AVX2-vectorized path; ACC has a single hand-written path.

All three implementations make the same three changes to ML-KEM:
\begin{enumerate}
  \item the module rank $k$ is scaled to $\{7,18,24\}$, giving the parameter
        sets kaiburr-1792, -4608, and -6144 (named kaiburr4, kaiburr6, and
        kaiburr8 in the code, after their noise parameter $t$; see
        Table~\ref{tab:kaiburr-params});
  \item ciphertext compression is removed, so both ciphertext polynomials are
        serialized at full $12$-bit precision; and
  \item the centered binomial sampler is replaced by a sampler for the noise
        family $f(t)$, with $t\in\{4,6,8\}$.
\end{enumerate}
Everything else, including the ring, the modulus $q=3329$, the NTT and its constants,
the Fujisaki--Okamoto transform, and the symmetric primitives, is left
unchanged.

These modifications nonetheless require more than just mechanical changes to the code. ML-KEM
implementations assume $k \le 4$ in several places, and the AVX2 packers
assume a particular coefficient order. Most of the rewritten code follows from
scaling $k$, which we describe first (Section~\ref{sec:impl:k}). We then cover
the noise sampler (Section~\ref{sec:impl:sampler}), the removal of compression
(Section~\ref{sec:impl:compression}), and what is checked about each
implementation (Section~\ref{sec:impl:assurance}).

\subsection{Scaling $k$}
\label{sec:impl:k}

\paragraph{Accumulator overflow.}
Larger values of $k$ can overflow the matrix-vector product. In the NTT domain, each polynomial of the 
product is a sum of $k$ pointwise products, which
ML-KEM accumulates in an \texttt{int16\_t} and reduces once at the end. In the C reference, the 
C AVX2, and the Jasmin reference code, each output coefficient is the sum of two separately
Montgomery-reduced products, so it is bounded by $2(q-1) = 6656$ rather than by
$q$, and the sum of $k$ terms fits only while $2kq < 2^{15}$, i.e.\ $k \le 4$. On the other hand, the Jasmin AVX2 code accumulates 
both products at $32$ bits and applies a single Montgomery reduction, so its terms are bounded by $q$ and it is safe to
$k \le 9$; it also fails at $k=18$ and $k=24$.

The C and Jasmin implementations therefore Barrett-reduce the accumulator after every pointwise product, at 
a cost of one extra reduction per term.\footnote{We note that reducing after every term is more conservative than
necessary, since a Barrett reduction leaves headroom for several further terms.
We did not implement this, as reduction is not a significant part of the cost.}
ACC needs no such change since it accumulates with \texttt{bn.addvm}, 
a modular addition against the \texttt{MOD} register, which is set to $2q$ throughout the basemul. Its accumulator
stays below $2q$ for every $k$, and the reduction is free.

\paragraph{Matrix generation.}
The AVX2 matrix sampler in ML-KEM hardcodes how the $k^2$
matrix entries are grouped into batches of four for the $4$-way Keccak for each $k \le 4$. We
replace this with a generic loop; it samples four entries at a time with 4-way Keccak and handles 
any remaining entries with a scalar tail. Only $k=7$ has a
remainder, of one entry.

The same loop generates both $\mathbf{A}$ for key generation and its transpose
$\mathbf{A}^T$ for encapsulation, by selecting the matching domain separators.
The C code computes these separators from the row and column indices, while
the Jasmin code reads them from a precomputed table. ACC generates one entry at a time and 
derives each separator from a running nonce, with separate handling for the transposed and 
non-transposed orders.

\paragraph{Memory.}
The C and Jasmin code store the full matrix $\mathbf{A}$ on the stack. Its size
grows quadratically with $k$, from $24.5$\,KiB at $k=7$ to $288$\,KiB at
$k=24$. This is not suitable for ACC, whose
stock configuration only provides $32$\,KiB of data memory. That is less than the
working set of a single decapsulation at $k=24$, which is about $34$\,KiB for
the ciphertext, re-encrypted ciphertext, secret key, and hashing buffers alone,
before the matrix is counted.

We address this in two ways. First, the ACC implementation never stores
$\mathbf{A}$. It generates one polynomial at a time, multiplies it into a
running accumulator, and discards it. Second, we enlarge ACC's visible data
memory from $31$\,KiB to $127$\,KiB in the hardware description, out of
$128$\,KiB of physical DMEM including a $1$\,KiB scratch region. The ACC
results in Section~\ref{sec:benchmarks} are therefore for this modified hardware
configuration.

\subsection{The noise sampler}
\label{sec:impl:sampler}

The noise family $f(t)$ maps $t$ uniformly random bits $a_1,\ldots,a_t$, read
least-significant first, to the coefficient
$$(2a_t-1)\,\Bigl(1-a_1+2\prod_{i<t}a_i\Bigr).$$
All our implementations replace the centered binomial sampler with this
expression, and the PRF output grows to $t\cdot n/8$ bytes per polynomial. The
five variants differ in how they compute two parts of the expression: the
product term $\prod_{i<t}a_i$ and the sign $(2a_t-1)$.
Table~\ref{tab:sampler-variants} summarizes their choices.

\begin{table}[h]
  \centering
  \caption{How each implementation computes the product term and applies the
           sign in $f(t)$.}
  \label{tab:sampler-variants}
  \begin{tabular}{lllc}
    \toprule
    Variant & Product term & Sign & Word size \\
    \midrule
    C ref       & bit-by-bit AND chain   & integer multiply             & scalar \\
    C AVX2      & compare against a mask & conditional lane negate      & 8--32-bit lanes \\
    Jasmin ref  & carry propagation      & conditional two's-complement & scalar \\
    Jasmin AVX2 & compare against a mask & conditional two's-complement & 8--16-bit lanes \\
    ACC         & carry propagation      & modular subtract + select    & 16-bit lanes \\
    \bottomrule
  \end{tabular}
\end{table}

The Jasmin reference code and the ACC assembly avoid the literal AND chain
with a carry trick. Adding one to the low $t-1$ bits carries into bit $t-1$
exactly when all of those bits are set, so a right shift gives us the product.
The AVX2 kernels instead isolate the same bits and apply a vector compare,
which gives the product directly as an all-ones or all-zero mask. In every
vectorized and assembly variant, we then apply the sign separately and
without branching.

Extracting $t$-bit fields from the PRF output is simplest for $t=8$ and $t=4$,
which align with bytes and nibbles, and hardest for $t=6$, as explained in Section~\ref{sec:distribution}. Both AVX2 kernels use an extra 
shuffle or gather for this step. The C kernel also reads slightly beyond the end of the PRF output, so the buffer includes 32 bytes of 
trailing slack. ACC lacks a byte-shuffle instruction, so it instead uses a sequence of \texttt{bn.rshi} shifts to 
move the 6-bit fields across register boundaries.

Key generation needs $2k$ noise polynomials. The AVX2 code samples them in
groups of four using the $4$-way Keccak.

\subsection{Removing compression}
\label{sec:impl:compression}

\textsc{Kaiburr} serializes both ciphertext polynomials $\mathbf{u}$ and $v$
at full $12$-bit precision, so the ciphertext grows to $(k+1)\cdot 384$ bytes.
The public key keeps its size of $384k+32$ bytes, since ML-KEM never
compresses it.

The reference code and ACC reuse their existing $12$-bit packer unchanged. The
AVX2 code needs a new one. Its existing packer is used only for the public key,
which is stored in the permuted coefficient order of the AVX2 NTT, so it
applies that permutation while packing. Ciphertext polynomials leave the
inverse NTT in standard order. Both AVX2 forks therefore add a packer and
unpacker that work in standard order, without the permutation.

The new packer emits $24$ bytes per iteration as two $16$-byte stores twelve
bytes apart, so its final iteration would write past the end of the buffer.
The C code stages the final iteration through a scratch buffer, and the Jasmin
code handles it as a separate case.

Removing compression also has a consequence beyond engineering. It makes the decryption-failure probability
formally provable, which we develop in Section~\ref{sec:verification}.

\subsection{Assurance}
\label{sec:impl:assurance}

\paragraph{Jasmin.}
The Jasmin implementation is a fork of \texttt{formosa-mlkem}~\cite{fvk},
whose EasyCrypt proofs establish functional correctness for ML-KEM. These
proofs do not carry over to \textsc{Kaiburr}, since the sampler, module rank,
and serialization all change. What we check is much narrower.

Every exported function carries two annotations. A \texttt{ct} annotation is
checked by Jasmin's constant-time checker, for both standard and speculative
constant-time. A \texttt{safety} annotation establishes that array accesses
stay in bounds and that all memory is initialized before use. Both properties
are checked for all six builds.

\paragraph{Testing.}
The C implementation relies on the deterministic round-trip self-tests
inherited from its upstream ML-KEM codebase. The Jasmin implementation runs
the same kind of self-test, and additionally cross-checks its output against
the C implementation, bit-for-bit, across all six builds. The ACC
implementation is validated against a Python model based on
\texttt{kyber-py}\footnote{\url{https://github.com/GiacomoPope/kyber-py}} which subclasses \texttt{ML\_KEM}, overrides the noise sampler, and
disables ciphertext compression. For each parameter set, the model generates
known-answer tests (KATs) for key generation, encapsulation, and decapsulation
(including implicit rejection), as well as round-trip tests.
Unlike the Jasmin code, ACC's constant-time behavior is not machine-checked;
the hand-written kernels avoid secret-dependent branches
by construction.